\section*{Data Availability}
\label{sec:data}
This is a named submission (confirmed with the course instructor,
2026-11-13), so the replication package below is the real, public
repository, not an anonymized mirror:
\url{https://github.com/PASS-IITG/cs5187-term-project-dauntlesscrux}.

The replication package contains \tool's source
(\texttt{src/tsck/\allowbreak\{cfg.py, outparam.py, compare.py\}}) and its 19
engine unit tests; the rule as data for Protocol B
(\texttt{specs/outparam\_reset.json}) and Protocol A
(\texttt{specs/bio\_chain\_pop.json}, \texttt{protocol\_a.py}, with its
own regression tests in \texttt{test\_protocol\_a\_logic.py}); the
pinned OpenSSL commit SHAs for every snapshot in
Table~\ref{tab:rq1-wholetree} (\texttt{run\_scan.sh},
\texttt{run\_protocol\_a.sh}); and the unedited logs, JSON results, and an
\texttt{environment.txt} recording tool versions behind this paper's
numbers (\texttt{src/tsck/results/}), including an independent timing
rerun's log.

It also contains the manual triage labels and evidence behind every claim
in Section~\ref{sec:eval}: \texttt{triage.csv}/\texttt{triage.md} for
Protocol B; \texttt{protocol\_a\_triage.md} for Protocol A, with source
excerpts for all 11 whole-tree sites;
\texttt{reachability\_rq1\_candidates.md} for the corrected reachability
derivation behind the 2 potential-recurrence sites;
\texttt{unclear\_sites\_recheck.md} for the caller read that reclassifies
the 2 formerly-unclear \texttt{OSSL\_CMP\_get1\_rootCaKeyUpdate} sites as
convention deviation; and
\texttt{historical\_recurrence\_search.md} for 10 independently-authored
OpenSSL fixes of the same release-without-reset shape found by a
commit-history search. Finally, it includes the tool-comparison evidence
behind Table~\ref{tab:rq2-comparison}. \texttt{make check} reruns the
offline checks; \texttt{make scan} reruns every OpenSSL scan from a
clean checkout.

\section*{Author Contributions}
\todo{This is a template, not a statement of fact -- neither author has
filled it in yet. Replace each bracketed line with what that person
actually did, in enough detail that it would hold up if either of you
were asked about the other's part in the final presentation (per the
course rubric, "that was X's part" is not an acceptable answer). This
needs to be true, specific, and agreed by both of you before
submission, not generated by an assistant.}
Rutvij Wagh (230101113): [fill in]. Udayraj Choudhary (230101107):
[fill in]. Both authors reviewed and take responsibility for the full
paper and replication package, including the sections produced with
generative-AI assistance disclosed below.

\section*{Use of Generative AI}
\todo{Confirm and finalize before submission. This draft reflects what
one assistant was asked to do in one work session. Check it against what
both authors actually used across the whole project. Only then treat it
as a true statement of record.}

A generative AI assistant (\todo{name and version}) was used to draft
prose in several sections from material the authors had already
produced and verified: experiment logs, triage spreadsheets, and tool
output under \texttt{src/tsck/results/} and \texttt{research/}. It also
searched for and verified the five related-work citations in
Section~\ref{sec:related} against publisher pages and DOIs. For
Protocol~B, it did not run OpenSSL builds, design the rule, or produce
any underlying experimental data: all Protocol~B numbers were produced
by the authors running \tool and the comparison scripts.

Protocol~A (Section~\ref{sec:rq5}) is an exception needing closer
verification. The assistant designed its rule and
\texttt{specs/bio\_chain\_pop.json} from the authors' existing
CVE-2023-0215 recurrence notes; ran it directly against a real,
independently-cloned copy of the pinned OpenSSL source in its own
sandboxed environment (\texttt{results/protocol\_a\_run\_environment.txt}
discloses a slightly different \texttt{tree-sitter} version than the
project's prior pin); and wrote the false-positive classification for
all 11 whole-tree findings in \texttt{results/protocol\_a\_triage.md}
after reading the real source at the pinned commit. Authors reproduced
that run on a separate machine (Section~\ref{sec:setup}), but it was the
assistant who performed the first read of the flagged sites and wrote
the reasoning. The same caveat applies to three further pieces of primary
analysis the assistant produced and wrote up itself, not from
author-supplied material. The first is the corrected reachability
argument for the 2 potential-recurrence sites in
\texttt{results/reachability\_rq1\_candidates.md}, which also corrects a
factual error in the authors' own original triage note. The second is the
reclassification of the 2 previously-unclear
\texttt{OSSL\_CMP\_get1\_rootCaKeyUpdate} sites in
\texttt{results/unclear\_sites\_recheck.md}, performed by reading the
function's one real in-tree caller, \texttt{apps/cmp.c}, for the first
time. The third is the commit-history search behind
\texttt{results/historical\_recurrence\_search.md}. \todo{Authors must
independently re-check the Reason~A/B/C classification for each of the
11 Protocol A sites, the reachability argument for all 4
potential-recurrence/formerly-unclear sites, and a sample of the ten
historical commits cited, each against its own source before submission,
and confirm here, as fact not assumption, that they did -- in all four
cases the assistant produced the primary
analysis, not just prose describing already-verified work.}
